\section*{Apêndice A --- Planejamento das visualizações longitudinais}
\label{ap:visualizacoes}

As Tabelas~\ref{tab:planejamento-figuras} e~\ref{tab:planejamento-figuras-cont}
apresentam o planejamento das visualizações que serão geradas durante a
execução do TCC~II, relacionando cada figura prevista à respectiva questão
de pesquisa, à métrica principal e ao tipo de visualização adotado. O
conjunto foi desenhado para permitir a comparação entre as quatro épocas
históricas E1 (2009--2013), E2 (2014--2017), E3 (2018--2021) e E4
(2022--atual), conforme delimitação operacional na
Seção~\ref{sec:metodologia}, e para evidenciar \emph{(i)}~tendências
evolutivas de longo prazo (E1~$\rightarrow$~E3) e \emph{(ii)}~o eventual
impacto da Era da IA assistida (E3~$\rightarrow$~E4).

\begin{table}[htbp]
\centering
\caption{Planejamento das visualizações por questão de pesquisa (parte~1).}
\label{tab:planejamento-figuras}
\scriptsize
\setlength{\tabcolsep}{3pt}
\renewcommand{\arraystretch}{1.05}
\begin{tabularx}{\linewidth}{|c|c|>{\raggedright\arraybackslash}p{0.20\linewidth}|>{\raggedright\arraybackslash}p{0.18\linewidth}|>{\raggedright\arraybackslash}X|}
\hline
\textbf{Fig.} & \textbf{Q} & \textbf{Métrica} & \textbf{Visualização} & \textbf{Descrição} \\
\hline
F1 & Q1 & CBO &
\emph{Boxplot} &
Acoplamento médio entre classes por época (mediana, dispersão e
\emph{outliers}). \\
\hline
F2 & Q1 & RFC &
\emph{Boxplot} &
Métodos potencialmente acionados por classe em cada época. \\
\hline
F3 & Q1 & CBO, RFC &
Série temporal &
Medianas de CBO e RFC ao longo das quatro épocas. \\
\hline
F4 & Q2 & WMC &
\emph{Boxplot} &
Complexidade lógica agregada por classe em cada época. \\
\hline
F5 & Q2 & \texttt{avg\_loc\_method} &
\emph{Boxplot} &
Tamanho médio dos métodos; hipótese de métodos maiores na Era da IA. \\
\hline
\end{tabularx}
\end{table}

\begin{table}[htbp]
\centering
\caption{Planejamento das visualizações por questão de pesquisa (parte~2).}
\label{tab:planejamento-figuras-cont}
\scriptsize
\setlength{\tabcolsep}{3pt}
\renewcommand{\arraystretch}{1.05}
\begin{tabularx}{\linewidth}{|c|c|>{\raggedright\arraybackslash}p{0.20\linewidth}|>{\raggedright\arraybackslash}p{0.18\linewidth}|>{\raggedright\arraybackslash}X|}
\hline
\textbf{Fig.} & \textbf{Q} & \textbf{Métrica} & \textbf{Visualização} & \textbf{Descrição} \\
\hline
F6 & Q2 & \texttt{avg\_wmc}, \texttt{avg\_loc\_method} &
Dispersão &
Correlação entre tamanho e complexidade médios, colorida por época. \\
\hline
F7 & Q3 & DIT &
\emph{Violin plot} &
Profundidade média da hierarquia de herança por época. \\
\hline
F8 & Q3 & NOC &
\emph{Violin plot} &
Número médio de subclasses diretas por classe. \\
\hline
F9 & Q3 & LCOM\_HS &
\emph{Boxplot} &
Ausência média de coesão por repositório em cada época. \\
\hline
F10 & Q4 & Todas &
Painel comparativo E3 vs.\ E4 &
Contraste da Era da IA, com \emph{p-valores} ajustados (Mann--Whitney). \\
\hline
\end{tabularx}
\end{table}

{\footnotesize Elaborado pelos autores. Os artefatos finais (PNG,
\emph{notebooks} Python/R e arquivo \texttt{.pbix} do Power~BI) serão
publicados no pacote de replicação previsto para o TCC~II.}

\medskip
Para garantir a comparabilidade entre figuras, todas as visualizações
adotarão a mesma paleta de cores por época, escala consistente
entre painéis e identificação explícita das épocas E1 (2009--2013),
E2 (2014--2017), E3 (2018--2021) e E4 (2022--atual). Sempre que
cabível, os \emph{p-valores} ajustados (Bonferroni) dos testes de
Kruskal--Wallis e Mann--Whitney (E3 \emph{vs.} E4) serão reportados no
próprio painel ou em legenda associada.
